% ---- V. Results (P2) -- target ~2 pages ----
% PREFER report/results_auto.tex (python -m src.write_results), which fills every number in automatically.
% This hand-fill template is kept as a fallback.
% Tables and figures come from: python -m src.metrics && python -m src.plots
% On Overleaf, upload results/tables/*.tex to tables/ and results/figures/*.pdf to figures/.
% Fill every \tbd{} from results/key_numbers.md. Do not type numbers that are not in that file.
\providecommand{\tbd}[1]{\textcolor{red}{[#1]}}

\section{Results}
\label{sec:results}

All percentages are over 150 pairs (300 calls per experiment and model). Brackets give 95\% Wilson
confidence intervals, and parse failures are scored as wrong. Where two conditions are run on the same
calls, we compare them with an exact McNemar test.

\subsection{Pairwise judging and position bias (Exp.~1)}

\input{tables/t1_pairwise.tex}

\begin{figure}[t]
\centering
\includegraphics[width=\columnwidth]{figures/f1_position_bias.pdf}
\caption{Pairwise accuracy when the correct answer is shown first (A) or second (B). The dashed line is
chance. Error bars are 95\% Wilson CIs.}
\label{fig:position}
\end{figure}

Table~\ref{tab:t1_pairwise} and Fig.~\ref{fig:position} show that both judges are \tbd{more / less}
accurate when the correct answer appears as \tbd{A}. For the 3B model, accuracy is \tbd{x}\% with the correct
answer in position A and \tbd{y}\% in position B, a gap of \tbd{g} percentage points. For the 7B model the gap is
\tbd{g}. The 3B judge outputs \texttt{A} in \tbd{p}\% of its parsed verdicts, against the 50\% that an unbiased
judge would produce on this balanced design. Only \tbd{c}\% (3B) and \tbd{c}\% (7B) of pairs receive the same pick in
both orders. In the remaining pairs the verdict follows the position and not the code: \tbd{AA} pairs got
\texttt{A} twice and \tbd{BB} got \texttt{B} twice.
\tbd{One sentence: does the larger model reduce position bias, or only raise accuracy?}

\subsection{Pointwise judging (Exp.~3)}

\input{tables/t2_pointwise.tex}

\begin{figure}[t]
\centering
\includegraphics[width=\columnwidth]{figures/f2_pointwise.pdf}
\caption{Pointwise judging of all 300 answers. FAR: incorrect code accepted. FRR: correct code rejected.}
\label{fig:pointwise}
\end{figure}

Judged one answer at a time, the 3B model reaches \tbd{x}\% accuracy and the 7B model \tbd{y}\%
(Table~\ref{tab:t2_pointwise}). Accuracy alone hides the failure that matters for code auditing.
The false acceptance rate is \tbd{FAR}\% for 3B and \tbd{FAR}\% for 7B, so \tbd{about k in 10} incorrect solutions
would pass review. The false rejection rate is only \tbd{FRR}\%. The judges label \tbd{p}\% of all answers
\texttt{correct}, which shows a leniency bias \tbd{(confirm: is it well above 50\%?)}. Only \tbd{q}\% of pairs
have both their answers judged right.

\subsection{Effect of the fixes (Exp.~4)}

\input{tables/t3_fixes.tex}

\begin{figure*}[t]
\centering
\includegraphics[width=0.85\textwidth]{figures/f3_fixes.pdf}
\caption{Accuracy (left) and order consistency (right) for the baseline and the three fixes. Order averaging
is consistent by construction and is scored only on the pairs where both orders agree.}
\label{fig:fixes}
\end{figure*}

\textbf{Order averaging (4a)} uses no extra model calls beyond Exp.~1. It keeps a verdict only when both orders agree.
On those pairs, accuracy rises to \tbd{x}\% (3B) and \tbd{y}\% (7B), but only \tbd{c}\% and \tbd{c}\% of pairs get a
verdict. The rest are left undecided. \textbf{The debias prompt (4b)} moves the share of \texttt{A} verdicts from \tbd{p}\% to
\tbd{q}\% (3B), and changes accuracy by \tbd{d} points (McNemar $p=\tbd{}$). \textbf{Majority voting (4c)}
over three samples at $T{=}0.7$ \tbd{does / does not} improve on greedy decoding (\tbd{x}\% vs.\ \tbd{y}\%,
$p=\tbd{}$). This suggests the bias is \tbd{systematic, not sampling noise}.

\subsection{Verbosity (Exp.~2)}

\input{tables/t4_verbosity.tex}

\begin{figure}[t]
\centering
\includegraphics[width=\columnwidth]{figures/f4_verbosity.pdf}
\caption{How often the 3B judge picks the incorrect answer before and after that answer is padded with comments
claiming correctness. The code logic is unchanged (checked by AST equality).}
\label{fig:verbosity}
\end{figure}

Adding comments to the incorrect answer, with its logic unchanged, changes how often it wins from
\tbd{x}\% to \tbd{y}\% (\tbd{+d} points, McNemar $p=\tbd{}$). \tbd{Interpretation: the judge is / is not swayed by
comments that claim correctness.} Because padding also makes the incorrect answer the longer one in
127 of 150 pairs (Section~\ref{sec:data}), this effect combines persuasive comments with extra length.

\subsection{By difficulty}

\input{tables/t5_difficulty.tex}

\tbd{One or two sentences on whether accuracy or consistency falls on hard problems. Add the caveat that
CodeJudgeBench difficulty labels disagree across generator splits for 82\% of the problems that appear in more than
one split, so this breakdown is indicative only.}

\subsection{Summary of findings}
\begin{enumerate}
\item \tbd{Finding 1: position bias (size of the gap, consistency).}
\item \tbd{Finding 2: pointwise FAR and what it means for auditing.}
\item \tbd{Finding 3: which fix works, and its cost (coverage or extra calls).}
\item \tbd{Finding 4: 3B vs 7B.}
\item \tbd{Finding 5 (if Exp.~2 ran): verbosity.}
\end{enumerate}
